\documentclass[11pt,a4paper]{article}
\usepackage[margin=2.4cm]{geometry}
\usepackage[T1]{fontenc}
\usepackage{lmodern}
\usepackage{amsmath,amssymb}
\usepackage{graphicx}
\usepackage{booktabs}
\usepackage{adjustbox}
\usepackage{xcolor}
\usepackage[numbers,sort&compress]{natbib}
\usepackage{caption}
\usepackage{enumitem}
\usepackage[hidelinks]{hyperref}
\usepackage{cleveref}

\captionsetup{font=small,labelfont=bf}
\setlist{itemsep=2pt,topsep=4pt}
\newcommand{\pendingbox}[1]{\fbox{\parbox{0.9\linewidth}{\centering\small\textit{#1}}}}
\newcommand{\fig}[3]{%
  \IfFileExists{#1}{\includegraphics[width=#2]{#1}}{\pendingbox{Figure pending: #3}}}
\newcommand{\todo}[1]{\par\noindent\colorbox{yellow!25}{\parbox{\dimexpr\linewidth-2\fboxsep}{\small\textbf{To update when the run finishes:} #1}}\par}

\newcommand{\ResvisionMedian}{361\,m}
\newcommand{\ResvisionmaggravMedian}{341\,m}
\newcommand{\ResvisionmaggravValMedian}{94.0\,m}
\newcommand{\ResSSFixRate}{71\%}
\newcommand{\ResSSFixMedian}{1.0\,m}
\newcommand{\ResPFFullMedian}{1.4\,m}
\newcommand{\ResPFFullWithin}{98\%}
\newcommand{\ResPFBestMedian}{1.4\,m}
\newcommand{\ResPFBestWithin}{98\%}
\newcommand{\ResvisionValMedian}{79.8\,m}
\newcommand{\ResvisionmagMedian}{370\,m}
\newcommand{\ResvisionmagValMedian}{95.2\,m}
\newcommand{\RescoarseMedian}{136\,m}
\newcommand{\RescoarseValMedian}{102\,m}
\newcommand{\RescoarsebalancedMedian}{145\,m}
\newcommand{\RescoarsebalancedValMedian}{105\,m}


\title{\textbf{Enabling Onboard Navigation for UAVs in GPS-Denied Environments}\\[0.4em]
\large Multimodal localisation from aerial imagery, magnetometry and gravimetry\\
with uncertainty-aware fusion and coarse-to-fine refinement}
\author{Advika Sinha (23117010) \and Cherish Puniani (23117042) \and Namay Rohatgi (23119022) \and Varun R Mallya (23117144)\\[0.6em]
\textit{Advisor:} Prof.\ Nitish Arya\\
Indian Institute of Technology Roorkee\\
MIP400A: B.Tech Project}
\date{\today}

\begin{document}
\maketitle

\begin{abstract}
A UAV that loses GNSS must localise itself from what its own sensors observe. We build and evaluate,
in simulation on real survey data, the multimodal localiser proposed in our mid-term evaluation: a
Vision Transformer for the aerial image and 1D-CNN encoders for magnetic and gravity profiles, each
producing a probabilistic representation $(\mu,\Sigma)$, fused by an uncertainty-aware core (UDML)
that computes noise-aware uncertainty $\rho$, modality dependency $\alpha$ and fusion weights
$w_v,w_m,w_g$, followed by a localisation head. We assemble a co-registered dataset over Mountain Pass,
California from four years of 0.6\,m NAIP imagery and a 2016 helicopter survey with a full-tensor
gravity gradiometer and a caesium magnetometer, and evaluate on held-out flight lines (whose readings
were never gridded into the training maps) with held-out-year imagery. Regressing coordinates
directly, the model reaches a median error of \ResvisionmaggravValMedian{} on training-year imagery but
\ResvisionmaggravMedian{} on the held-out year. We therefore wrap it in a coarse-to-fine pipeline: the
network proposes candidate map cells, magnetic/gravity profile matching (TERCOM) selects among them,
and a learned image matcher (LoFTR) against an onboard reference map refines the position. Reporting
a position only when the image match is confident, a single frame yields a fix in \ResSSFixRate{} of
cases with a median error of \ResSSFixMedian{} and no wrong fixes. Tracking over whole flights with a
particle filter that combines all three modalities, weighted by UDML, gives a median error of
\ResPFFullMedian{} (\ResPFFullWithin{} of steps within 10\,m). We report ablations per modality, the
behaviour of the UDML weights, and the limitations of the simulation.
\end{abstract}

\tableofcontents
\clearpage

%======================================================================
\section{Introduction}
GNSS signals are weak and easily jammed or spoofed, yet survey, disaster-response and defence UAVs must
keep flying when GNSS fails. Classical solutions predate GNSS: TERCOM matched radar-altimeter terrain
profiles against elevation maps~\cite{golden1980tercom} and DSMAC correlated downward-looking images with
a reference scene~\cite{carr1980dsmac}. Modern counterparts match UAV images to satellite imagery with
learned features~\cite{goforth2019gps,patel2020visual,hierarchical2025} or match magnetic and gravity
anomaly measurements against geophysical maps~\cite{canciani2016absolute,gravityreview2026}.

This semester's goal was an end-to-end pipeline in simulation: aerial observations in, latitude and
longitude out. Our contributions are:
\begin{enumerate}
\item A co-registered multimodal dataset (imagery, magnetometry, gravimetry, positions) over a
250\,km$^2$ region, with an evaluation protocol that holds out entire flight lines from the
geophysical maps and an entire imagery year (\cref{sec:data}).
\item An implementation of the mid-term architecture (\cref{fig:arch}): ViT and 1D-CNN encoders with
probabilistic outputs and the UDML fusion core (\cref{sec:model}).
\item A coarse-to-fine system around it, combining learned coarse localisation, retrieval, TERCOM-style
profile matching, learned image matching and particle-filter tracking (\cref{sec:c2f,sec:pf}), and
literature-driven fixes for the weaknesses we measured (\cref{sec:phase3}).
\item A quantitative evaluation including per-modality ablations and the behaviour of the UDML weights
under sensor degradation (\cref{sec:results}).
\end{enumerate}

\begin{figure}[t]
\centering
\fig{figures/architecture.png}{0.95\linewidth}{architecture}
\caption{The architecture proposed at the mid-term evaluation and implemented here.}
\label{fig:arch}
\end{figure}

%======================================================================
\section{Related work}
\paragraph{Visual absolute localisation.}
Goforth and Lucey~\cite{goforth2019gps} align UAV frames to satellite imagery with CNN features trained
for cross-season robustness and report an average error below 8\,m on a 0.85\,km flight. Patel et
al.~\cite{patel2020visual} localise against Google Earth renderings with sub-metre RMSE on a 303\,m
flight. SWA-PF~\cite{swapf2025} and Shan et al.~\cite{shan2017googlemap} combine map matching with
particle filters; FoundLoc~\cite{foundloc2023} uses foundation-model features without an initial pose.
Recent work splits the problem into retrieval and fine matching~\cite{hierarchical2025,anyvisloc2025}.
Benchmarks with real UAV imagery include UAV-VisLoc~\cite{uavvisloc2024} and AnyVisLoc~\cite{anyvisloc2025}.

\paragraph{Retrieval and matching.}
DINOv2~\cite{oquab2023dinov2} features give strong zero-shot place recognition (AnyLoc~\cite{keetha2023anyloc});
contrastive training with hard negatives (Sample4Geo~\cite{deuser2023sample4geo}, TransGeo~\cite{zhu2022transgeo})
is the state of the art for drone-to-satellite retrieval. LoFTR~\cite{sun2021loftr} and
LightGlue~\cite{lindenberger2023lightglue} with DISK~\cite{tyszkiewicz2020disk} provide dense and sparse
learned correspondences. Direct pose regression~\cite{kendall2015posenet} is known to behave like image
retrieval and to generalise poorly~\cite{sattler2019limits}.

\paragraph{Geophysical navigation.}
Canciani and Raquet~\cite{canciani2016absolute} report about 13\,m horizontal error over an hour-long
flight with a caesium magnetometer and a marginalised particle filter, a best case at low altitude.
Gravity matching is reviewed in~\cite{gravityreview2026}. MagNav~\cite{gnadt2020magnav,magnavjl}
provides reference filters.

\paragraph{Multimodal learning and sequences.}
A dominant modality can starve the others of gradient; OGM-GE~\cite{peng2022ogm} modulates gradients to
balance them. SeqSLAM~\cite{milford2012seqslam} and SeqNet~\cite{garg2021seqnet} match sequences rather than
single frames to cope with appearance change.

%======================================================================
\section{Data}\label{sec:data}
\paragraph{Site.} Mountain Pass, Mojave Desert, California (lon $-115.58$ to $-115.38$, lat $35.36$ to
$35.59$). Terrain elevation ranges from 867 to 2399\,m; the area contains mountains, alluvial fans, a
dry lake, solar plants, a mine and the I-15 corridor (\cref{fig:data}).

\paragraph{Imagery.} NAIP~\cite{naip} 4-band orthoimagery at 0.6\,m for 2016, 2018, 2020 and 2022
(81 quarter-quad tiles, 36.5\,GB), mosaicked per year and resampled to 1.25\,m/px. A 224\,px input crop
covers about 280\,m.

\paragraph{Geophysics.} A USGS-contracted helicopter survey (late 2016)~\cite{usgs_mtnpass_agg,usgs_mtnpass_mag}
with a HeliFALCON airborne gravity gradiometer and a Scintrex CS-3 caesium magnetometer: 184 flight
lines, traverse lines 200\,m apart (100\,m infill in part) and tie lines, median 68\,m above ground,
gravity sampled at 8\,Hz (about 3.4\,m along-track) and magnetics decimated to about 27\,m. We use the
vertical gravity $g_D$ and the gradient $G_{DD}$ \emph{without} terrain correction (what an aircraft
senses) and the residual magnetic intensity (RMI). Everything is expressed in UTM 11N (NAD83).

\paragraph{Splits.} Traverse lines are ordered across-track; of every eight adjacent lines one is
used for validation and one for testing, and each infill line follows its neighbour; tie lines stay in
training. The magnetic and gravity \emph{training maps are gridded only from training lines} (linear
interpolation at 25\,m with light smoothing; samples within 60\,m of a held-out line are dropped), so
evaluation never uses readings that went into the maps. Validation and test samples are the raw
readings of the held-out lines, resampled to 25\,m steps and cut into 64-step (1.6\,km) windows:
1{,}597 validation and 3{,}063 test windows (21 and 20 lines). The 20 test segments, with windows every
50\,m, form the flights used for tracking. Training imagery is 2016/2018/2020; test imagery is 2022; the
onboard reference map for matching is 2020.

\begin{figure}[t]
\centering
\begin{minipage}{0.3\linewidth}\centering
\fig{figures/generated/naip_2016.jpg}{\linewidth}{NAIP mosaic}\end{minipage}\hfill
\begin{minipage}{0.68\linewidth}\centering
\fig{figures/generated/geophysics.png}{\linewidth}{geophysics grids}\end{minipage}
\caption{Left: the 2016 NAIP mosaic of the study area. Right: magnetic RMI, $g_D$ and $G_{DD}$ grids
built from training lines only, with validation (green) and test (black) lines overlaid.}
\label{fig:data}
\end{figure}

\paragraph{Training samples.} Training samples are generated on the fly at random continuous positions:
a synthetic 64-step track with random heading and slight curvature is sampled from the training maps,
with sensor noise, offsets and gain errors (magnetics $\sigma=2$\,nT, offset 20\,nT; $g_D$ 0.3\,mGal;
$G_{DD}$ 5\,E), and an image crop from a random training year with rotation ($\pm10^\circ$), scale
($\pm15\%$), colour and blur augmentation. Each modality is dropped with probability 0.15.

%======================================================================
\section{The mid-term model}\label{sec:model}
\paragraph{Encoders.} The image encoder is a DINOv2 ViT-S/14~\cite{dosovitskiy2021vit,oquab2023dinov2},
fine-tuned with layer-wise learning-rate decay; its CLS token is mapped to $\mu_v$ and $\log\sigma^2_v$
($d=256$). The magnetic and gravity encoders are 1D-CNNs (four convolutional blocks, global mean/max
pooling, MLP) over the profile, its first difference, the height above ground and the heading.

\paragraph{UDML core.} For modalities $k\in\{v,m,g\}$ with diagonal Gaussians $(\mu_k,\sigma_k^2)$:
\begin{align}
\rho_k &= \tfrac{1}{d}\textstyle\sum_i \sigma_{k,i}^2 &&\text{(noise-aware uncertainty)}\\
\alpha_k &= \tfrac{1}{|\mathcal{A}|-1}\textstyle\sum_{j\in\mathcal{A}, j\neq k} A_{jk} &&\text{(dependency: attention $k$ receives)}\\
w_k &= \operatorname{softmax}_{k\in\mathcal{A}}\left(\log\alpha_k - \log\rho_k\right) &&\text{(fusion weights)}
\end{align}
where $A$ is the attention matrix of one multi-head attention layer over the three modality tokens and
$\mathcal{A}$ the available modalities. The fused feature is $z=\sum_k w_k \hat z_k$ with
$\hat z_k=\mu_k+\sigma_k\epsilon$ during training and $\hat z_k=\mu_k$ at test time.

\paragraph{Heads and loss.} The regression head outputs normalised easting/northing and an isotropic
log-variance, trained with a Gaussian negative log-likelihood. The classification head outputs logits
over 250\,m map cells, trained with soft labels (Gaussian over neighbouring cells, $\sigma=0.5$ cell).
The same head is applied to each $\mu_k$ alone (auxiliary outputs, weight 0.3) to report per-modality
accuracy and to keep every encoder trained; a small KL term ($\beta=10^{-4}$) regularises the latents.
Training uses AdamW, bf16, batch 64, 30 epochs of 50k samples on one RTX 5070 Ti; the three
regression stages follow the mid-term timeline (vision, $+$magnetometry, $+$gravimetry).

%======================================================================
\section{Coarse-to-fine localisation}\label{sec:c2f}
\paragraph{Coarse.} The network's cell probabilities (classifier) or a Gaussian around its regressed
position (regressor) form a prior over 250\,m cells; optionally multiplied by a retrieval prior.

\paragraph{Retrieval.} The 2020 reference map is cut into 224\,px tiles every 100\,m and embedded by
the slide-6 ViT ([CLS, mean patch] $\to$ MLP $\to$ unit vector). We evaluate zero-shot DINOv2 features
(AnyLoc-style~\cite{keetha2023anyloc}) and contrastive fine-tuning (Sample4Geo-style~\cite{deuser2023sample4geo}):
symmetric InfoNCE between augmented training-year queries and reference tiles, with half of the
batches drawn from one 2\,km neighbourhood (hard negatives) and pairs closer than 250\,m excluded as
negatives. The ViT is initialised from the trained classifier.

\paragraph{Profile matching (TERCOM).} For every position on a 25\,m lattice inside the top-10 cells
(plus 125\,m margin) and headings $\{-3^\circ,0,3^\circ\}$, the measured 1.6\,km magnetic and gravity
profiles are compared with the training maps along the dead-reckoned track (1\% drift):
\begin{equation}
\log L(p) = -\frac{1}{2\tau}\sum_{c}\sum_{i=1}^{64}\frac{\big(r_{c,i}-\bar r_c\big)^2}{\sigma_c^2},\qquad
r_{c,i}=y_{c,i}-M_c(p+\Delta_i),
\end{equation}
removing each channel's unknown offset $\bar r_c$. The noise levels are calibrated on validation lines
as the residual between real readings and the map at the true track ($\sigma_\text{mag}=12.6$\,nT,
$\sigma_{g_D}=0.28$\,mGal, $\sigma_{G_{DD}}=37$\,E); $\tau=4$ accounts for correlated samples. The
posterior combines $\log L$ with the coarse prior.

\paragraph{Image matching (DSMAC).} The UAV view is simulated from the test-year mosaic with heading
($\sigma=8^\circ$), scale ($\sigma=7\%$) and camera-tilt (up to $\pm10^\circ$, homography keeping the
nadir point at the image centre) errors. Around the five best profile-matching hypotheses it is matched
against the 2020 reference map within 250\,m, with LoFTR~\cite{sun2021loftr} (kornia~\cite{riba2020kornia})
and a RANSAC similarity transform. A position is reported only if the match has at least 15 inliers;
otherwise the output is \emph{no fix} and the vehicle continues on its inertial solution.

%======================================================================
\section{Particle-filter tracking}\label{sec:pf}
Each held-out segment is flown as a sequence of steps 50\,m apart with 3000 particles~\cite{gordon1993pf,thrun2005probabilistic}:
\begin{itemize}
\item \textbf{Motion:} the dead-reckoned displacement with a per-flight heading/scale bias and random drift.
\item \textbf{Geophysics:} the profile likelihood above over the in-map part of the track, tempered for
the overlap between consecutive windows and scaled by $g_k=\mathrm{clip}(3w_k,0.25,2)$ from the UDML weights.
\item \textbf{Image:} LoFTR around the main particle clusters; a confident fix adds a robust Gaussian
likelihood ($\sigma=5$\,m, outlier floor 0.02, scaled by $g_v$) and seeds 10\% of the particles at the fix.
\item \textbf{Recovery:} a robust coarse-prior term $\log(1+C\,p_\text{cell})$ that rewards but never
punishes, and 2\% of particles re-drawn from the prior each step.
\item \textbf{Estimate:} the weighted mean of the mode holding most weight; systematic resampling when the
effective sample size drops below $N/2$.
\end{itemize}

\section{Literature-driven fixes}\label{sec:phase3}
\begin{itemize}
\item \textbf{Balanced training.} OGM~\cite{peng2022ogm}: each step, a modality whose auxiliary loss is
lower than the others' mean by $\delta$ has its encoder gradients scaled by $1-\tanh(0.1\,\delta)$;
vision is additionally dropped 40\% of the time.
\item \textbf{Sequence prior.} SeqSLAM-style~\cite{milford2012seqslam,garg2021seqnet}: the coarse
log-priors of the last five frames are averaged after shifting by the dead-reckoned displacement.
\item \textbf{Colour alignment.} Each query crop is given the reference map's per-channel mean and
standard deviation (only image statistics, so available onboard).
\item \textbf{Dead-reckoning bias estimation.} Each particle carries a heading and scale bias that
corrects its motion (the augmented-state idea of marginalised-PF magnetic navigation~\cite{canciani2016absolute,magnavjl}).
\end{itemize}

%======================================================================
\section{Results}\label{sec:results}
All numbers are on the held-out test lines with 2022 imagery unless stated otherwise.

\subsection{The mid-term model on its own}
\begin{table}[t]
\centering
\small
\caption{The slide-6 model on its own: median localisation error of the fused output on the validation lines (training-year imagery) and the test lines (held-out 2022 imagery), the error of each modality's auxiliary head alone, and the mean UDML weights.}
\label{tab:baseline}
\begin{adjustbox}{max width=\linewidth}
\begin{tabular}{lrrrrrrc}
\toprule
Model & Val.\ (2016--20) & Test (2022) & Test p90 & Vision alone & Mag alone & Grav alone & $w_v / w_m / w_g$ \\
\midrule
Vision only & 79.8\,m & 361\,m & 4{,}560\,m & 361\,m & -- & -- & 1.00 / 0.00 / 0.00 \\
Vision + magnetometer & 95.2\,m & 370\,m & 4{,}266\,m & 376\,m & 2{,}386\,m & -- & 0.97 / 0.03 / 0.00 \\
Vision + mag + gravity & 94.0\,m & 341\,m & 3{,}590\,m & 358\,m & 3{,}123\,m & 2{,}503\,m & 0.94 / 0.02 / 0.03 \\
Cell classifier (V+M+G) & 102\,m & 136\,m & 1{,}642\,m & 139\,m & 572\,m & 633\,m & 0.93 / 0.04 / 0.03 \\
Balanced classifier (V+M+G) & 105\,m & 145\,m & 3{,}823\,m & 152\,m & 585\,m & 655\,m & 0.91 / 0.05 / 0.04 \\
\bottomrule
\end{tabular}
\end{adjustbox}
\end{table}

\begin{table}[t]
\centering
\small
\caption{Mean UDML weights $w_v / w_m / w_g$ when one input is corrupted at test time (image: $\sigma=1$, profiles: $\sigma=2$ in normalised units). A working uncertainty-aware fusion lowers the weight of the corrupted modality.}
\label{tab:udml}
\begin{adjustbox}{max width=\linewidth}
\begin{tabular}{lcccc}
\toprule
Model & Clean & Noisy image & Noisy mag & Noisy grav \\
\midrule
Vision + magnetometer & 0.97 / 0.03 / 0.00 & 0.78 / 0.22 / 0.00 & 0.92 / 0.08 / 0.00 & -- \\
Vision + mag + gravity & 0.94 / 0.02 / 0.03 & 0.83 / 0.07 / 0.11 & 0.93 / 0.04 / 0.03 & 0.95 / 0.02 / 0.03 \\
\bottomrule
\end{tabular}
\end{adjustbox}
\end{table}


\Cref{tab:baseline} shows the central limitation of direct regression: on validation lines with
training-year imagery the full model reaches \ResvisionmaggravValMedian{}, but with held-out 2022
imagery the median rises to \ResvisionmaggravMedian{}, with a heavy tail. This is consistent with the
known behaviour of pose regression, which interpolates between memorised appearances rather than
measuring position~\cite{sattler2019limits}. The magnetic and gravity encoders on their own localise
only to kilometres: a 1.6\,km profile along one direction is ambiguous, and the neural encoders do not
learn the map-matching that explicit profile correlation performs.

\Cref{tab:udml} shows that the UDML weights concentrate on vision (about 0.95). When the image is
corrupted, weight moves to the geophysical modalities, as intended; when a profile is corrupted, its
weight does not consistently decrease. In the trained models, UDML therefore behaves mostly as
``trust vision'', which motivated the balanced training of \cref{sec:phase3}.

\subsection{Coarse-to-fine, one frame at a time}
\begin{table}[t]
\centering
\small
\caption{Single-shot coarse-to-fine pipeline on the 3{,}063 test windows: median error after each stage, and the image matcher's error when given a hypothesis at the true position.}
\label{tab:c2fstages}
\begin{adjustbox}{max width=\linewidth}
\begin{tabular}{lrrrrrr}
\toprule
Prior & Right cell in top-10 & Coarse & + TERCOM mag & + TERCOM grav & + TERCOM both & LoFTR at truth \\
\midrule
Vision regressor & 56\% & 372\,m & 141\,m & 135\,m & 97.0\,m & 1.0\,m \\
V+M regressor & 55\% & 391\,m & 142\,m & 135\,m & 97.4\,m & 1.6\,m \\
V+M+G regressor & 57\% & 360\,m & 140\,m & 126\,m & 93.8\,m & 1.6\,m \\
Cell classifier & 89\% & 136\,m & 77.5\,m & 72.2\,m & 58.3\,m & 1.6\,m \\
Cell classifier, vision prior & 86\% & 139\,m & 80.6\,m & 75.1\,m & 60.8\,m & 1.6\,m \\
\bottomrule
\end{tabular}
\end{adjustbox}
\end{table}

\begin{table}[t]
\centering
\small
\caption{Single-shot output with the \emph{no-fix} policy: a position is reported only when the LoFTR match has at least 15 RANSAC inliers. Wrong fixes are reported positions more than 25\,m off, as a share of all frames.}
\label{tab:c2ffix}
\begin{adjustbox}{max width=\linewidth}
\begin{tabular}{lrrrrr}
\toprule
Prior & Fix rate & Wrong fixes & Median of fixes & Fixes $<$10\,m & p90 of fixes \\
\midrule
Vision regressor & 71\% & 0.00\% & 1.0\,m & 100\% & 2.1\,m \\
V+M regressor & 70\% & 0.00\% & 1.5\,m & 100\% & 2.9\,m \\
V+M+G regressor & 70\% & 0.00\% & 1.6\,m & 100\% & 3.0\,m \\
\bottomrule
\end{tabular}
\end{adjustbox}
\end{table}


Each stage of \cref{tab:c2fstages} narrows the position: from a few hundred metres (coarse) to about
100\,m (TERCOM with both profiles; each profile alone gives 130--140\,m) to about a metre (LoFTR).
Given a hypothesis at the truth, LoFTR is within 10\,m for 95--97\% of the windows despite the year
difference and the simulated camera errors. The limiting factor is the coarse step: when the true
cell is not among the candidates, no later stage recovers. \Cref{tab:c2ffix} shows that the matcher's
confidence identifies these cases: every reported fix is within 10\,m and wrong fixes are absent; the
remaining frames are reported as \emph{no fix}.

\begin{table}[t]
\centering
\small
\caption{Image matchers searching $\pm$250\,m around a hypothesis 100\,m from the truth (300 validation windows, 2022 query vs 2020 map). Realistic views add $\sigma=8^\circ$ heading, $\sigma=7\%$ scale and up to $\pm10^\circ$ camera tilt.}
\label{tab:matchers}
\begin{adjustbox}{max width=\linewidth}
\begin{tabular}{llrrrr}
\toprule
Views & Matcher & Median & $<$5\,m & $<$10\,m & ms/match \\
\midrule
\multicolumn{6}{c}{\textit{pending}} \\
\bottomrule
\end{tabular}
\end{adjustbox}
\end{table}

\begin{table}[t]
\centering
\small
\caption{Retrieval against the embedded 2020 reference map (tiles every 100\,m). A query counts as a hit when a retrieved tile lies within 250\,m of the truth.}
\label{tab:retrieval}
\begin{adjustbox}{max width=\linewidth}
\begin{tabular}{llrrrr}
\toprule
Embedder & Queries & R@1 & R@5 & R@10 & R@25 \\
\midrule
Zero-shot DINOv2 ViT-S/14 & val 2018 & 28\% & 44\% & 53\% & 64\% \\
Zero-shot DINOv2 ViT-S/14 & test 2022 & 27\% & 45\% & 55\% & 67\% \\
Zero-shot DINOv2 ViT-B/14 & val 2018 & 23\% & 38\% & 46\% & 57\% \\
Zero-shot DINOv2 ViT-B/14 & test 2022 & 22\% & 40\% & 48\% & 61\% \\
Fine-tuned slide-6 ViT (InfoNCE) & val 2018 & 100\% & 100\% & 100\% & 100\% \\
Fine-tuned slide-6 ViT (InfoNCE) & test 2022 & 79\% & 88\% & 91\% & 94\% \\
\bottomrule
\end{tabular}
\end{adjustbox}
\end{table}

\begin{table}[t]
\centering
\small
\caption{Single-shot variants on the test windows (realistic views unless noted).}
\label{tab:singleshot}
\begin{adjustbox}{max width=\linewidth}
\begin{tabular}{lrrrrr}
\toprule
Configuration & Right cell in top-$k$ & Fix rate & Wrong fixes & Median of fixes & Fixes $<$10\,m \\
\midrule
Classifier, NCC & 89\% & 21\% & 1.01\% & 4.7\,m & 84\% \\
Classifier, LoFTR & 89\% & 87\% & 0.00\% & 1.6\,m & 100\% \\
Classifier $\times$ retrieval, LoFTR & 93\% & 88\% & 0.00\% & 1.6\,m & 100\% \\
\quad same, easy views & 93\% & 90\% & 0.00\% & 1.0\,m & 100\% \\
\quad + sequence prior (5 frames) & 93\% & 88\% & 0.00\% & 1.6\,m & 100\% \\
\quad + colour alignment & 92\% & 89\% & 0.00\% & 1.6\,m & 100\% \\
\quad + top-30 cells & 97\% & 90\% & 0.00\% & 1.6\,m & 100\% \\
\quad + all three & 94\% & 88\% & 0.00\% & 1.6\,m & 100\% \\
Balanced classifier + all three & 93\% & 88\% & 0.00\% & 1.6\,m & 100\% \\
No slide-6 model: retrieval only & 89\% & 87\% & 0.00\% & 1.6\,m & 100\% \\
\bottomrule
\end{tabular}
\end{adjustbox}
\end{table}


\todo{interpret \cref{tab:matchers,tab:retrieval,tab:singleshot} (phase 2 and 3 single-shot results).}

\subsection{Tracking over flights}
\begin{table}[t]
\centering
\small
\caption{Particle-filter tracking over the 20 held-out test flights (3{,}063 steps, 50\,m apart), realistic views.}
\label{tab:track}
\begin{adjustbox}{max width=\linewidth}
\begin{tabular}{lrrrrrr}
\toprule
Configuration & Median & $<$10\,m & $<$25\,m & p90 & $<$10\,m after 10 steps & Lost flights \\
\midrule
\multicolumn{7}{c}{\textit{pending}} \\
\bottomrule
\end{tabular}
\end{adjustbox}
\end{table}


\begin{figure}[t]
\centering
\fig{figures/generated/track_full.png}{0.85\linewidth}{particle-filter track (phase 2)}
\caption{A held-out test flight: true track (green) and particle-filter estimate (red) over the 2022
imagery, and the error per step with image fixes marked.}
\label{fig:track}
\end{figure}

\todo{interpret \cref{tab:track} and \cref{fig:track}: overall accuracy, lost flights, the contribution of
each modality, UDML versus equal weights, and the phase-3 fixes.}

\subsection{Comparison with published results}
Published systems report, on real flights, below 8\,m average over 0.85\,km~\cite{goforth2019gps},
below 1\,m RMSE over 303\,m~\cite{patel2020visual}, below 10\,m~\cite{swapf2025}, about 13\,m for
magnetic-anomaly navigation~\cite{canciani2016absolute}, and within about 20\,m without an initial
pose~\cite{foundloc2023}. Our numbers are in the same range or better, but the comparison is not
like-for-like: our query images are simulated from orthophotos of the same programme as the
reference map, our dead reckoning is simulated, and we evaluate one 250\,km$^2$ region. A fair claim
requires the evaluation on real-flight benchmarks proposed in \cref{sec:future}.

%======================================================================
\section{Limitations}
\begin{itemize}
\item \textbf{Simulated camera.} Queries are warped orthophoto crops; real imagery adds motion blur,
lens distortion, relief displacement and different illumination.
\item \textbf{Simulated motion.} Dead reckoning is the true displacement plus modelled bias and drift,
not an IMU/visual-odometry solution.
\item \textbf{One region and one survey.} All results are from Mountain Pass; the geophysical maps exist
only where surveys were flown, and the gravity gradiometer used is not small-UAV-mountable.
\item \textbf{Absolute accuracy.} The reference imagery and the survey positions are themselves accurate
to about a metre, so sub-metre absolute accuracy cannot be demonstrated with this data.
\end{itemize}

\section{Future work}\label{sec:future}
\begin{itemize}
\item Evaluate the vision chain on real UAV footage: Goforth and Lucey's dataset~\cite{goforth2019gps},
UAV-VisLoc~\cite{uavvisloc2024} and AnyVisLoc~\cite{anyvisloc2025}; a brute-force LoFTR baseline over the
whole map.
\item Replace simulated dead reckoning with an IMU model (bias, random walk) and visual odometry, and
measure accuracy across image-matching outages.
\item A second surveyed region, and a public release of the co-registered dataset.
\item Fast, low-power deployment on an FPGA (next semester).
\end{itemize}

\section{Conclusion}
The mid-term architecture is implemented and trained end to end on a real multimodal survey. On its
own, direct coordinate regression does not generalise across imagery years, and the UDML weights
collapse onto vision. Used as the coarse stage of a coarse-to-fine system, together with profile
matching of real magnetic and gravity readings and learned image matching against an onboard map, the
same model yields confident single-frame fixes at metre level with no wrong fixes, and continuous
tracking with a particle filter. The open problems are the realism of the simulation and making the
multimodal fusion itself, not only the refinement, benefit from the geophysical sensors.

\bibliographystyle{unsrtnat}
\bibliography{references}

\end{document}
